\documentclass[
  12pt,
  oneside,
  a4paper,
  brazil
]{abntex2}

% --- Codificação e língua ---
\usepackage[T1]{fontenc}
\usepackage[utf8]{inputenc}
\usepackage{lmodern}
\usepackage[brazil]{babel}

% --- Pacotes essenciais ---
\usepackage{graphicx}
\usepackage{indentfirst}
\usepackage{microtype}
\usepackage{xcolor}

% --- Pseudocódigo e blocos de código ---
\usepackage{listings}
\lstset{
  basicstyle=\small\ttfamily,
  frame=single,
  breaklines=false,
  keepspaces=true,
  columns=fixed,
  backgroundcolor=\color[gray]{0.97},
  xleftmargin=1em,
  xrightmargin=1em,
}

% --- Citações ABNT ---
\usepackage[alf,abnt-etal-cite=2,abnt-year-extra-label=yes,abnt-etal-list=0]{abntex2cite}

% --- Metadados do documento ---
\titulo{Searchable Encryption: Busca em Dados Cifrados}
\autor{Artur Luiz Rizzato Toru Soda \and Davi Menegaz Junkes \and Matheus Fernandes Bigolin}
\local{Florianópolis}
\data{2026}
\instituicao{%
  Universidade Federal de Santa Catarina \par
  Centro Tecnológico \par
  Departamento de Informática e Estatística%
}
\tipotrabalho{Trabalho de disciplina}
\preambulo{%
  Trabalho apresentado à disciplina INE5429: Segurança em Computação
  do Curso de Ciências da Computação da Universidade Federal de Santa
  Catarina, como requisito parcial para a obtenção de nota na avaliação
  final. Professores: Jean Everson Martina e Thaís Bardini Idalino.%
}

\makeindex

\begin{document}

\frenchspacing

\imprimircapa
\imprimirfolhaderosto

\pdfbookmark[0]{\contentsname}{toc}
\tableofcontents*
\cleardoublepage

% ============================================================
\chapter{Introdução}
% ============================================================

Nos últimos anos, a computação em nuvem tornou-se a forma mais comum
de armazenar e compartilhar dados. A conveniência é evidente: redução
de custos, acesso de qualquer lugar e escalabilidade imediata. O
problema é que, ao enviar dados para um servidor remoto, o usuário
perde o controle sobre eles. Um provedor desonesto, ou um servidor
comprometido por um atacante, pode acessar, copiar ou vender essas
informações sem que o usuário perceba \cite{wang2016}.

A solução mais imediata é cifrar os dados antes de enviá-los. Assim,
mesmo que o servidor acesse os arquivos, só verá sequências de bits
sem significado, pois as chaves de decifração ficam com o usuário. O
problema é que a cifragem elimina a capacidade de busca: palavras
cifradas não podem ser comparadas diretamente, e encontrar um documento
específico exigiria baixar todo o conteúdo, decifrá-lo localmente e
só então realizar a busca. Essa abordagem é impraticável para qualquer
volume considerável de dados.

Searchable Encryption (SE) resolve essa tensão. Trata-se de um
conjunto de técnicas criptográficas que permite realizar buscas
diretamente sobre dados cifrados, sem revelar o conteúdo dos
documentos nem a palavra que está sendo buscada. O servidor executa a
busca e devolve os resultados corretos, mas aprende o mínimo possível
no processo.

O trabalho implementa o esquema proposto por Song, Wagner e Perrig em
\textit{Practical Techniques for Searches on Encrypted Data}
\cite{song2000}, o primeiro da área com provas formais de segurança. A
escolha se justifica pela clareza conceitual: as operações são simples
o suficiente para que a lógica de segurança possa ser acompanhada
passo a passo, ao mesmo tempo em que o esquema exibe limitações
concretas que motivaram décadas de pesquisa posterior. O restante do
trabalho apresenta a teoria do esquema, um exemplo com valores reais,
detalhes de implementação e três experimentos que verificam corretude,
desempenho e o vazamento de informação característico do esquema.

% ============================================================
\chapter{Detalhamento Teórico}
% ============================================================

\section{Modelo}

O esquema envolve três partes. O \textit{dono dos dados} guarda as
chaves criptográficas, cifra os documentos antes de enviá-los e é o
único capaz de gerar tokens de busca. O \textit{servidor} armazena os
documentos cifrados e executa buscas quando recebe um token, mas
nunca tem acesso às chaves. \textit{Usuários autorizados} podem
solicitar tokens ao dono para realizar consultas sem precisar ver o
conteúdo.

O modelo de ameaça adotado é o do servidor honesto mas curioso
(\textit{honest-but-curious}): o servidor segue o protocolo
corretamente, mas registra e analisa tudo o que recebe na tentativa de
extrair informações adicionais. Não se assume que ele vai modificar
resultados ou recusar buscas; apenas que vai tentar aprender o máximo
que puder a partir dos dados que passam por suas mãos. Esse modelo
captura o risco real de serviços de nuvem que, embora confiáveis
operacionalmente, podem ser curiosos quanto ao conteúdo armazenado.

\section{O Esquema de Song, Wagner e Perrig (2000)}

Para que a busca em dados cifrados funcione, o esquema precisa cumprir
dois objetivos ao mesmo tempo. Primeiro, a mesma palavra deve parecer
diferente em cada posição do documento, para que o servidor não
consiga identificar padrões só de olhar para o ciphertext. Segundo,
um único token de busca deve ser capaz de encontrar todas as
ocorrências dessa palavra, independentemente de quantas vezes ela
aparece. Esses dois objetivos parecem contraditórios à primeira vista,
mas a construção do esquema os resolve de forma elegante.

A ideia central é separar a cifragem em duas camadas. A primeira é
determinística: a função $E$, com chave $k''$, sempre produz o mesmo
resultado para a mesma palavra. Isso significa que $E(k'', \text{``the''})$
é idêntico em todas as posições onde ``the'' aparece, e esse valor é
chamado de $X_i$. É justamente essa consistência que permite ao token
de busca encontrar todas as ocorrências. A segunda camada existe para
ocultar $X_i$ do servidor. Para cada posição $i$, o esquema gera um
valor aleatório $S_i$ diferente e o usa para cobrir $X_i$ com uma
operação XOR. O que o servidor armazena, portanto, é apenas o
resultado desse XOR, sem ter acesso a $X_i$ diretamente.

Com essa intuição estabelecida, o algoritmo de cifragem pode ser
descrito com mais precisão. O conjunto de chaves é gerado aleatoriamente
e compõe três valores: $k_{\mathit{seed}}$ (semente do gerador de
sequência), $k'$ (chave de derivação) e $k''$ (chave da cifra
determinística). Para cifrar cada palavra $W_i$, primeiro calcula-se
$X_i = E(k'', W_i)$, que tem $n$ bits. Esse valor é dividido ao meio:
$L_i$ são os primeiros $n{-}m$ bits e $R_i$ são os últimos $m$ bits.
A partir de $L_i$, deriva-se uma chave específica para a posição:
$k_i = f(k', L_i)$. Em paralelo, o gerador de sequência produz $S_i$
a partir de $k_{\mathit{seed}}$ indexado pela posição $i$. Com $k_i$ e
$S_i$, calcula-se um verificador $P_i = F(k_i, S_i)$. O bloco de
cobertura é $T_i = S_i \| P_i$, com o mesmo tamanho de $X_i$, e o
ciphertext final é $C_i = X_i \oplus T_i$.

Para buscar uma palavra $W$, o usuário calcula $X = E(k'', W)$ e
$k = f(k', L)$, onde $L$ é a metade esquerda de $X$, e envia o par
$(X, k)$ ao servidor. Esse par é o trapdoor: contém exatamente o que
o servidor precisa para realizar a busca, mas não revela a palavra
buscada, pois inverter a cifra $E$ sem a chave $k''$ é
computacionalmente inviável.

O servidor recebe $(X, k)$ e percorre cada bloco $C_i$. Para cada um,
calcula $V = C_i \oplus X$ e verifica se $F(k, V_{\mathit{esq}})$ é
igual a $V_{\mathit{dir}}$, onde $V_{\mathit{esq}}$ e $V_{\mathit{dir}}$
são as metades esquerda e direita de $V$. Quando a palavra na posição
$i$ é a palavra buscada, $X_i = X$, e o XOR cancela: $V$ resulta em
$T_i = S_i \| P_i$. Como $k = k_i$ nesse caso, a verificação
$F(k, S_i) = P_i$ passa. Quando a palavra é diferente, o XOR não
cancela, $V$ não tem estrutura e a verificação falha com probabilidade
$1 - 1/2^m$. Com $m = 64$, essa probabilidade é praticamente 1,
tornando os falsos positivos desprezíveis.

A decifração não envolve o servidor. O usuário regenera a sequência
$S_i$ a partir de $k_{\mathit{seed}}$ (o gerador é determinístico),
desfaz o XOR para recuperar $X_i$ e aplica $E^{-1}(k'', X_i)$ para
obter a palavra original.

\section{Exemplo Concreto}

Para tornar o funcionamento do esquema mais concreto, os algoritmos
foram aplicados ao documento \texttt{["the", "cat", "sat", "on",
"the", "mat"]}, com $n = 128$ bits e $m = 64$ bits. Os valores abaixo
foram produzidos pela execução direta do código implementado e estão
representados em hexadecimal.

\textbf{KeyGen} gerou as seguintes chaves:

\begin{lstlisting}
k_seed = ebcaa1a50d8dda82 5592931d32ecf516
k'     = 6dfc6551bc6c8548 fb141965eaccfd3e
k''    = f638607a7f6d3df6 e12c3454d442f149
\end{lstlisting}

\textbf{Encrypt} para a palavra ``cat'' na posição 1. A palavra é
convertida para um bloco de 16 bytes com bytes nulos como padding e a
cifragem segue os passos descritos na Seção 2.2:

\begin{lstlisting}
W_1 = 6361740000000000 0000000000000000  ("cat" + 13 bytes 0x00)

X_1 = E(k'', W_1) = c72fe676b0f5cca5 47bc601746ac02d5
L_1 = X_1[:8]     = c72fe676b0f5cca5
R_1 = X_1[8:]     = 47bc601746ac02d5

S_1 = G(k_seed)[1] = 4a854df46c69f653
k_1 = f(k', L_1)   = 87b25aa5b62cec8a 596fc66d885c5786
P_1 = F(k_1, S_1)  = 9a3adc38b890a198
T_1 = S_1 || P_1   = 4a854df46c69f653 9a3adc38b890a198

C_1 = X_1 XOR T_1  = 8daaab82dc9c3af6 dd86bc2ffe3ca34d
\end{lstlisting}

Aplicando o mesmo processo a todas as seis palavras, obtêm-se os
seguintes blocos cifrados:

\begin{lstlisting}
C[0]  "the"  ->  eb6de9d830b6a889 59f9b4b52d90b86a
C[1]  "cat"  ->  8daaab82dc9c3af6 dd86bc2ffe3ca34d
C[2]  "sat"  ->  b1915b9b29933112 fe2a8562327f331d
C[3]  "on"   ->  11a16a04a238abb0 23bfc06d267c9b6b
C[4]  "the"  ->  170b562afc654a00 881511d34d313758
C[5]  "mat"  ->  ff74fda8fddf9d85 ff09a6f626a3c531
\end{lstlisting}

\textbf{Trapdoor} para buscar ``cat''. O usuário calcula
$X = E(k'', \text{``cat''})$ e $k = f(k', X[:8])$ e envia o par ao
servidor:

\begin{lstlisting}
X = c72fe676b0f5cca5 47bc601746ac02d5
k = 87b25aa5b62cec8a 596fc66d885c5786
\end{lstlisting}

\textbf{Search}: o servidor aplica $V = C_i \oplus X$ e verifica
$F(k, V[:8]) = V[8:]$ em cada bloco. Retorna a posição~1 para ``cat'',
as posições 0 e~4 para ``the'', e lista vazia para ``dog'', que não
existe no documento. Os valores gerados neste exemplo também serão
usados na próxima seção para ilustrar as propriedades de segurança.

\section{Propriedades de Segurança}

Duas propriedades de segurança do esquema ficam evidentes a partir do
exemplo anterior.

A primeira é o \textbf{sigilo provável}. A palavra ``the'' aparece nas
posições 0 e~4, mas os respectivos blocos cifrados são completamente
diferentes:

\begin{lstlisting}
C[0] = eb6de9d830b6a889 59f9b4b52d90b86a
C[4] = 170b562afc654a00 881511d34d313758
\end{lstlisting}

Isso acontece porque cada posição recebe um $S_i$ diferente gerado
pelo PRG. Como $C_i = X_i \oplus T_i$ e $T_i$ depende de $S_i$, o
resultado final é distinto mesmo que $X_i$ seja o mesmo nas duas
posições. O servidor não consegue saber, apenas olhando para o
ciphertext, que $C[0]$ e $C[4]$ correspondem à mesma palavra.

A segunda é a \textbf{consulta oculta}. O trapdoor enviado ao servidor
é o par $(X, k)$, onde $X = E(k'', \text{``cat''})$. Sem a chave $k''$,
não é possível recuperar ``cat'' a partir de $X$. O servidor consegue
apenas verificar se um determinado bloco casa com o trapdoor; a
palavra buscada, em si, permanece desconhecida.

Essas duas propriedades garantem privacidade do conteúdo dos
documentos e das consultas, mas o esquema não oferece uma terceira
propriedade que frequentemente se deseja: a ocultação do padrão de
acesso.

\section{Limitações}

Toda vez que o servidor executa uma busca, ele observa quais posições
são retornadas. Ao longo de muitas consultas, esse histórico acumula
informação: palavras frequentes no documento retornam mais posições a
cada busca, e palavras raras retornam menos. Um servidor
honesto-mas-curioso pode analisar essa distribuição e inferir, com boa
precisão, a frequência relativa das palavras no documento, sem decifrar
nenhum bloco. Isso é o que se chama de vazamento por padrão de acesso
(\textit{access pattern leakage}), e é demonstrado experimentalmente
na Seção~4.3.

O esquema também é estático: a sequência $S_1, \ldots, S_l$ é gerada
uma única vez a partir de $k_{\mathit{seed}}$ no momento da cifragem.
Adicionar ou remover palavras exigiria re-cifrar o documento inteiro.
Esquemas mais recentes, como o de Curtmola et al. (2006), citado em
\cite{wang2016}, resolvem esse problema usando índices cifrados que
suportam atualizações, mas com maior custo de armazenamento.

% ============================================================
\chapter{Implementação}
% ============================================================

\section{Primitivas e Decisões de Projeto}

O esquema foi implementado em Python~3 usando a biblioteca
\textit{cryptography}. Cada elemento abstrato do paper corresponde a
uma primitiva concreta, listada no Quadro~\ref{quadro:primitivas}.

\begin{quadro}[htb]
\caption{Mapeamento de primitivas criptográficas}
\label{quadro:primitivas}
\begin{tabular}{p{2.8cm}p{5.5cm}p{4.8cm}}
\hline
\textbf{Símbolo} & \textbf{Função no esquema} & \textbf{Primitiva utilizada} \\
\hline
$E$ / $E^{-1}$ & Cifra determinística das palavras & AES-128 modo ECB \\
$f$ & Derivação de $k_i$ a partir de $L_i$ & HMAC-SHA256, truncado a 16 bytes \\
$F$ & PRF de verificação no servidor & HMAC-SHA256, truncado a 8 bytes \\
$G$ (PRG) & Geração dos valores $S_i$ & AES-128 modo CTR com IV~=~0 \\
\hline
\end{tabular}
\fonte{Elaborado pelos autores.}
\end{quadro}

A escolha do AES-ECB para $E$ é intencional: o esquema requer que a
cifragem seja determinística, ou seja, que a mesma palavra sempre gere
o mesmo $X_i$, independentemente da posição. Sem isso, o trapdoor não
conseguiria encontrar todas as ocorrências de uma palavra, pois cada
$X_i$ seria diferente mesmo para palavras iguais. Outros modos, como
CBC ou GCM, introduzem aleatoriedade e não podem ser usados aqui por
essa razão.

O PRG foi implementado com AES-CTR e IV zero. Isso garante que a
sequência $S_1, \ldots, S_l$ seja reprodutível a partir de
$k_{\mathit{seed}}$: basta reiniciar o gerador com a mesma chave e o
mesmo IV para obter exatamente os mesmos valores, o que é necessário
para que a decifração funcione sem armazenar os valores $S_i$.

Os parâmetros foram definidos como $n = 128$ bits (16 bytes, alinhado
ao bloco do AES) e $m = 64$ bits. Palavras menores que 16 bytes
recebem bytes nulos como padding; palavras maiores são truncadas. A
taxa de falsos positivos na busca é $1/2^{64}$ por posição, o que na
prática é nula para qualquer tamanho de documento.

\section{Pseudocódigo}

Os Algoritmos~\ref{alg:keygen} a~\ref{alg:decrypt} apresentam o
pseudocódigo dos cinco algoritmos do esquema.

\begin{lstlisting}[caption={KeyGen},label=alg:keygen]
Entrada: parametro de seguranca
Saida:   K = (k_seed, k', k'')

1. k_seed  <- bits_aleatorios(n)
2. k'      <- bits_aleatorios(n)
3. k''     <- bits_aleatorios(n)
4. retornar K
\end{lstlisting}

\begin{lstlisting}[caption={Encrypt},label=alg:encrypt]
Entrada: K = (k_seed, k', k''),  D = (W_1, ..., W_l)
Saida:   C = (C_1, ..., C_l)

1. S[1..l] <- PRG(k_seed, l)         // l valores de (n-m) bits
2. para i = 1 ate l:
3.     X_i <- E(k'', W_i)            // cifra deterministica
4.     L_i <- X_i[1 .. n-m]
5.     k_i <- f(k', L_i)             // chave especifica da posicao
6.     T_i <- S[i] || F(k_i, S[i])  // bloco de cobertura
7.     C_i <- X_i XOR T_i
8. retornar C
\end{lstlisting}

\begin{lstlisting}[caption={Trapdoor},label=alg:trapdoor]
Entrada: K = (k_seed, k', k''),  W (palavra a buscar)
Saida:   T_W = (X, k)

1. X <- E(k'', W)
2. L <- X[1 .. n-m]
3. k <- f(k', L)
4. retornar T_W = (X, k)
\end{lstlisting}

\begin{lstlisting}[caption={Search (executado pelo servidor)},label=alg:search]
Entrada: C = (C_1, ..., C_l),  T_W = (X, k)
Saida:   lista de posicoes com correspondencia

1. posicoes <- lista vazia
2. para i = 1 ate l:
3.     V   <- C_i XOR X
4.     s   <- V[1 .. n-m]
5.     chk <- V[n-m+1 .. n]
6.     se F(k, s) = chk:
7.         adicionar i a posicoes
8. retornar posicoes
\end{lstlisting}

\begin{lstlisting}[caption={Decrypt},label=alg:decrypt]
Entrada: K = (k_seed, k', k''),  C = (C_1, ..., C_l)
Saida:   D = (W_1, ..., W_l)

1. S[1..l] <- PRG(k_seed, l)
2. para i = 1 ate l:
3.     L_i <- C_i[1..n-m]     XOR  S[i]
4.     k_i <- f(k', L_i)
5.     R_i <- C_i[n-m+1..n]   XOR  F(k_i, S[i])
6.     X_i <- L_i || R_i
7.     W_i <- E_inverso(k'', X_i)
8. retornar D
\end{lstlisting}

% ============================================================
\chapter{Resultados Experimentais}
% ============================================================

Os três experimentos a seguir verificam, respectivamente, a corretude
da implementação, o comportamento de desempenho previsto pela teoria e
a principal vulnerabilidade do esquema. Todos foram executados com o
código presente no repositório do projeto.

\section{Corretude}

O objetivo deste experimento é verificar que os cinco algoritmos
produzem resultados corretos em situações típicas de uso. Foi criado
um documento de 300 palavras escolhidas aleatoriamente de um
vocabulário fixo, inserindo a palavra ``target'' manualmente nas
posições 42, 150 e~275. Após a cifragem, cinco buscas foram
realizadas cobrindo os casos mais relevantes do esquema, cujos
resultados estão no Quadro~\ref{quadro:corretude}.

\begin{quadro}[htb]
\caption{Resultados do experimento de corretude}
\label{quadro:corretude}
\begin{tabular}{p{5.2cm}p{3.2cm}p{3.2cm}p{1.5cm}}
\hline
\textbf{Caso de teste} & \textbf{Esperado} & \textbf{Obtido} & \textbf{Status} \\
\hline
``target'' (3 ocorrências) & {[42, 150, 275]} & {[42, 150, 275]} & OK \\
``qwerty'' (inexistente)   & {[]}             & {[]}             & OK \\
``cloud'' (múltiplas)      & todas as pos.    & todas as pos.    & OK \\
Decifração (300 palavras)  & doc. original    & doc. original    & OK \\
Trapdoor com chave errada  & {[]}             & {[]}             & OK \\
\hline
\end{tabular}
\fonte{Elaborado pelos autores.}
\end{quadro}

O quinto caso verifica o isolamento entre chaves diferentes. Um segundo
conjunto de chaves $K_2$ foi gerado, o mesmo documento foi cifrado com
$K_2$, e tentou-se buscar nesse ciphertext usando um trapdoor gerado
com a chave original $K$. O resultado foi uma lista vazia, confirmando
que trapdoors de uma chave não produzem correspondências válidas em
ciphertexts cifrados com outra chave.

\section{Desempenho}

O segundo experimento verifica que o tempo de busca cresce linearmente
com o número de palavras, como previsto pela análise teórica de
complexidade $O(n)$ de \citeonline{song2000}. Para cada tamanho
testado, o documento foi cifrado e o tempo médio de cinco execuções da
busca foi medido. A busca foi sempre realizada por uma palavra ausente
do documento, de forma a forçar a varredura completa sem interrupção
antecipada, correspondendo ao pior caso.

\begin{figure}[htb]
\caption{Tempo médio de busca em função do número de palavras}
\label{fig:desempenho}
\begin{center}
\includegraphics[width=0.85\textwidth]{grafico_desempenho}
\end{center}
\fonte{Elaborado pelos autores.}
\end{figure}

Os resultados confirmam o comportamento linear: de 500 palavras (cerca
de 2~ms) a 50.000 palavras (cerca de 180~ms), o tempo cresce
proporcionalmente ao tamanho do documento. O custo por iteração é uma
operação XOR de 16 bytes seguida de um cálculo de HMAC-SHA256. Como
não há estrutura de índice, o servidor não tem como saber em quais
blocos está a palavra buscada sem varrer todos.

\section{Vazamento por Frequência de Acesso}

O terceiro experimento demonstra a limitação de segurança descrita na
Seção~2.5. O objetivo é mostrar que um servidor honesto-mas-curioso,
sem decifrar nenhum dado, consegue inferir a frequência relativa das
palavras no documento apenas observando os padrões de busca ao longo
do tempo.

Foi criado um vocabulário de sete palavras com frequências
propositalmente desiguais: \textit{cloud} (40 ocorrências),
\textit{data} (30), \textit{secure} (20), \textit{encrypt} (15),
\textit{search} (10), \textit{server} (6) e \textit{private} (2). O
documento foi cifrado e, em seguida, simulou-se um usuário realizando
300 buscas, escolhendo qual palavra buscar com probabilidade
proporcional à sua frequência real. O servidor registrou o número de
posições retornadas em cada consulta e acumulou esse total por palavra
ao longo de todas as 300 buscas.

\begin{figure}[htb]
\caption{Frequência real versus padrão de acesso observado pelo servidor}
\label{fig:vazamento}
\begin{center}
\includegraphics[width=0.9\textwidth]{grafico_vazamento}
\end{center}
\fonte{Elaborado pelos autores.}
\end{figure}

A Figura~\ref{fig:vazamento} mostra que o ranking inferido pelo
servidor a partir dos padrões de acesso reflete com precisão a
frequência real das palavras no documento. Palavras mais frequentes
retornam mais posições a cada busca e, ao longo de muitas consultas,
acumulam muito mais acessos no histórico do servidor. Isso revela a
estrutura de frequência do documento sem que o servidor tenha lido uma
única palavra. O resultado confirma o que \citeonline{song2000} já
reconheciam: o esquema é vulnerável a ataques estatísticos baseados em
frequência, e esse foi um dos principais motivadores para o
desenvolvimento de esquemas SSE mais robustos nas décadas seguintes.

% ============================================================
\chapter{Conclusão}
% ============================================================

O trabalho implementou o esquema de busca em dados cifrados de
\citeonline{song2000} em Python~3, usando AES-128~ECB como cifra
determinística, HMAC-SHA256 como função pseudo-aleatória e AES-128~CTR
como gerador de sequência. Os cinco algoritmos do esquema foram
implementados, verificados com um mini-exemplo concreto e avaliados em
três experimentos distintos.

Os experimentos confirmaram as propriedades centrais do esquema. A
corretude foi verificada em cinco cenários: busca de palavras com
ocorrências conhecidas, busca de palavra inexistente, múltiplas
ocorrências, decifração completa e isolamento entre chaves. O
desempenho confirmou a complexidade $O(n)$: o tempo de busca cresce
proporcionalmente ao número de palavras, sem atalhos possíveis pela
ausência de estrutura de índice. O experimento de vazamento demonstrou
que um servidor honesto-mas-curioso consegue recuperar a distribuição
de frequências do documento apenas observando os padrões de acesso ao
longo de 300 buscas.

O esquema de \citeonline{song2000} é um marco histórico por ter
estabelecido as bases formais da área, mas suas limitações são
concretas. A estaticidade, que não permite atualizar o documento sem
re-cifrar tudo, e o vazamento por padrão de acesso foram os problemas
que impulsionaram a pesquisa seguinte. Esquemas baseados em índices
cifrados, discutidos em \cite{wang2016}, resolvem ambas as questões,
porém com maior custo de armazenamento e processamento. O experimento
de vazamento desenvolvido aqui ilustra de forma direta por que esses
avanços foram necessários.

\noindent\textit{Declaração sobre uso de inteligência artificial:}
Ferramentas de inteligência artificial foram utilizadas como auxílio na
revisão e na formatação do texto. O conteúdo técnico, as decisões de
projeto, a implementação dos algoritmos e os experimentos realizados
são de autoria dos integrantes do grupo.

% ============================================================
\postextual
\bibliography{referencias}
% ============================================================

\end{document}
